\section{推理与行动的融合：ReAct}

\subsection{Reasoning 范式与 Acting 范式}

在 ReAct 出现之前，LLM 的应用分裂为两个独立的范式：

\begin{itemize}
    \item \textbf{Reasoning}（如 Chain-of-Thought）：LLM 内部逐步推理，增强 test-time compute，但无法获取外部知识或工具。
    \item \textbf{Acting}（如 RAG、tool use）：LLM 调用外部环境（搜索引擎、计算器、API），获取知识和计算能力，但缺乏推理以适应复杂局面。
\end{itemize}

两个范式各自解决 QA 的不同子问题（推理类问题 vs. 知识类问题 vs. 计算类问题），导致方法碎片化（piecemeal）。

\subsection{ReAct：推理 + 行动的统一框架}

\begin{importantbox}{ReAct 的核心思想}
让 LLM 交替生成\textbf{思考}（thought）和\textbf{动作}（action），观测（observation）由环境返回。只需写一个示例轨迹作为 prompt，展示如何思考和行动来解决任务，LLM 就能模仿这种模式。
\end{importantbox}

具体工作流程：
\begin{enumerate}
    \item LLM 生成一个 thought（如 ``我需要查找这三家公司的市值''）
    \item LLM 生成一个 action（如 \texttt{Google{[}Apple market cap{]}}）
    \item 环境返回 observation（搜索结果）
    \item 将 thought + action + observation 追加到上下文，LLM 继续生成下一步
    \item 重复直到任务完成
\end{enumerate}

\begin{knowledgebox}{推理与行动的双向增益}
\begin{itemize}
    \item \textbf{行动帮助推理}：从外部获取实时信息、执行计算
    \item \textbf{推理帮助行动}：根据当前情况规划策略、在异常时重新调整计划。例如，当搜索返回股价而非市值时，推理可以指导下一步查找股份数量来间接计算。
\end{itemize}
\end{knowledgebox}

\subsection{ReAct 的通用性}

ReAct 不局限于问答——任何可以转化为``文本游戏''的任务都可以使用 ReAct 框架：

\begin{itemize}
    \item 网页浏览、在线购物
    \item 视频游戏（通过 image captioning 将像素转为文本）
    \item 家庭任务模拟（ALFWorld）
    \item 软件工程（SWE-bench）
\end{itemize}

\begin{warningbox}{没有推理的 Acting Agent 的局限}
在 ALFWorld 游戏中，纯 acting agent（不生成 thought）在遇到异常时（如目标物体不在预期位置）会反复执行同一动作，无法适应。加入 thinking action 后，agent 可以反思当前状况并重新规划。
\end{warningbox}

\subsection{本章小结}

ReAct 的贡献不仅是一个具体方法，更是提出了 Reasoning Agent 的范式——将推理作为 Agent 的一种特殊``内部动作''。这种范式系统性地优于纯推理或纯行动。

